\documentclass[../main]{subfiles}

\begin{document}

\chapter{Segundo Principio de la Termodinámica: Ciclo de Carnot}

\section{Introducción al Segundo Principio de la Termodinámica}

El Segundo Principio de la Termodinámica establece que en cualquier proceso espontáneo, la entropía del universo siempre tiende a aumentar, lo que se traduce en una pérdida de la disponibilidad de energía útil. Uno de los conceptos fundamentales relacionados con este principio es el Ciclo de Carnot, que proporciona una comprensión profunda de la eficiencia de los motores térmicos.

\subsection{El Ciclo de Carnot}

El Ciclo de Carnot es un ciclo termodinámico idealizado que opera entre dos fuentes de temperatura, una caliente ($T_h$) y otra fría ($T_c$). Este ciclo consta de cuatro procesos reversibles: dos isotérmicos y dos adiabáticos. Su eficiencia es máxima y está determinada únicamente por las temperaturas de las fuentes de calor.

\info{Ciclo de Carnot}{
La eficiencia del Ciclo de Carnot está dada por la fórmula:

\begin{equation}
\eta = 1 - \frac{T_c}{T_h}
\end{equation}

Donde:
\begin{itemize}
    \item $\eta$: Eficiencia del ciclo.
    \item $T_c$: Temperatura de la fuente fría.
    \item $T_h$: Temperatura de la fuente caliente.
\end{itemize}
}
\ejemplo{Ejemplo}{Calcular la eficiencia del Ciclo de Carnot cuando la temperatura de la fuente caliente es de $500\,K$ y la temperatura de la fuente fría es de $300\,K$.

\textbf{Datos:}
\begin{itemize}
    \item $T_h = 500\,K$
    \item $T_c = 300\,K$
\end{itemize}

Sustituyendo en la fórmula de eficiencia:

\begin{equation}
\eta = 1 - \frac{300\,K}{500\,K} = 1 - 0.6 = 0.4
\end{equation}

Por lo tanto, la eficiencia del Ciclo de Carnot en este caso es del $40\%$.
}

\actividad{Responder el siguiente cuestionario}

\begin{enumerate}
\item ¿Cuál es la importancia del Segundo Principio de la Termodinámica en la ingeniería y la tecnología?
\item ¿Qué significan los términos "proceso reversible" y "proceso irreversible" en el contexto de la termodinámica?
\item ¿Por qué el Ciclo de Carnot se considera un ciclo ideal y cuál es su relevancia en la práctica?
\item ¿Cómo afecta la temperatura de las fuentes de calor a la eficiencia del Ciclo de Carnot?
\item ¿Cuál es la diferencia entre la eficiencia del Ciclo de Carnot y la eficiencia de un motor real?
\item Explique brevemente por qué ningún motor térmico puede tener una eficiencia del $100\%$ según el Segundo Principio de la Termodinámica.
\item ¿Qué es la entropía y cómo se relaciona con el Segundo Principio de la Termodinámica?
\item ¿Cuáles son las limitaciones prácticas del Ciclo de Carnot en la ingeniería de motores térmicos?
\item ¿Qué son los procesos isotérmicos y adiabáticos en el Ciclo de Carnot y cuál es su importancia?
\item ¿Qué aplicaciones prácticas tiene el conocimiento del Ciclo de Carnot en la ingeniería y la tecnología moderna?
\end{enumerate}
\actividad{Resolver los siguientes problemas}

\begin{enumerate}
\item Una máquina térmica opera según el Ciclo de Carnot entre dos fuentes de temperatura. La temperatura de la fuente caliente es de $600\,K$ y la temperatura de la fuente fría es de $300\,K$. Si la máquina absorbe $2000\,J$ de calor de la fuente caliente en cada ciclo, ¿cuánto trabajo realiza la máquina en cada ciclo? Suponga que el proceso es reversible y que no hay pérdidas de energía.
\item Se tiene un refrigerador que opera según el Ciclo de Carnot entre dos fuentes de temperatura. La temperatura de la fuente caliente es de $300\,K$ y la temperatura de la fuente fría es de $250\,K$. Si el refrigerador debe extraer $5000\,J$ de calor de la cámara fría en cada ciclo, ¿cuánto trabajo debe realizar el refrigerador en cada ciclo? Suponga que el proceso es reversible y que no hay pérdidas de energía.
\item Una planta de energía utiliza el vapor de agua como fluido de trabajo en un ciclo de potencia basado en el Ciclo de Carnot. La temperatura de la fuente caliente es de $500\,°C$ y la temperatura de la fuente fría es de $20\,°C$. Si la planta produce $10\,MW$ de potencia, ¿cuánto calor se absorbe de la fuente caliente en cada segundo? Suponga que el proceso es reversible y que no hay pérdidas de energía.
\item Se desea construir un motor térmico que opere según el Ciclo de Carnot entre dos fuentes de temperatura. Si se requiere que la eficiencia del motor sea del $50\%$ y la temperatura de la fuente caliente sea de $1000\,K$, ¿cuál debe ser la temperatura de la fuente fría? Suponga que el proceso es reversible y que no hay pérdidas de energía.
\item Un refrigerador opera según el Ciclo de Carnot entre dos fuentes de temperatura. La temperatura de la fuente caliente es de $350\,K$ y la temperatura de la fuente fría es de $250\,K$. Si el refrigerador debe extraer $1000\,J$ de calor de la cámara fría en cada ciclo, ¿cuánto calor se absorbe de la fuente caliente en cada ciclo? Suponga que el proceso es reversible y que no hay pérdidas de energía.
\end{enumerate}

\end{document}
